\section{What the new descent and paired estimates leave open}
\label{dsx-status}

The standard abc conjecture remains unproved and undisproved. The global
target is still one constant \(C_\varepsilon\), for each
\(\varepsilon>0\), independent of the positive primitive triple
\(a+b=c\), such that
\[
 \log c-(1+\varepsilon)\log\operatorname{rad}(abc)
       \le C_\varepsilon .
\]
A disproof would require a fixed positive \(\varepsilon\) and actual
primitive triples on which this difference is unbounded above.

\paragraph{A closed fixed branch and the varying-family problem.}
The ordinary two-descent proves the squared-unit quotient has rank zero
and exactly three rational points. Together with the other two unit
classes and the actual common-source reconstruction, it excludes the
primitive square--cube intersection stated above. The descent is on the
entire Jacobian, with all closed-point norms, finite-prime conditions,
real signatures and the rational-point argument explicitly retained.
Older checkpoint statements leaving this particular unit class open
are superseded by these exact conclusions.

The separate ramified analysis also excludes the complete specified
three-adic common-source covers. Its stronger integer consequence says
that a total cube second norm forces the first norm to be prime to
three, excluding the first shape \(M=3R^h\) whenever the total second
exponent is divisible by three. The numerical equations themselves
can be locally soluble: the exact global-unit extraction is essential.

This is not a classification of arbitrary primitive abc triples.
No argument has shown that all hypothetical abc counterexamples belong
to the excluded perfect-power family. Nonunit residuals and other
exponents require their own analysis, and there is still no proved
height bound with the needed uniform dependence as the curves vary.
Coverage and uniformity are separate obligations.

\paragraph{A paired estimate and the remaining signed tail.}
The shifted-resultant theorem supplies an actual common-depth lcm
bound across distinct primes. The all-owner theorem controls its
specified positive paired portion simultaneously, with an explicitly
bounded exceptional set. It does not prove that high-depth labels
have suitable partners. Nor does it control all excess above the
partner's depth, exceptional roots, the dependent-root complement,
or the required representation of arbitrary abc inputs.

The full signed sum retains the negative contributions from depths
one and two. A bound on positive certificate cost is not a lower
bound on that signed sum, and a density-one statement is not a
pointwise theorem. New certificates must meet their actual arithmetic
congruences and prove a saving after all costs and credits have been
included. The surviving parent routes remain open.

\paragraph{The formal boundary.}
The fifteen new Lean declarations verify the explicitly described
finite arithmetic at thirteen. They are distinct from the previous
twenty-four nonlinear-selector declarations. Their source was freshly
compiled and every declaration's axiom dependencies were inspected.
Neither count measures the proportion of an abc proof that has been
completed. The complete descent, local heights, and rational-point
classification remain ordinary mathematics under internal independent
review; they are not silently included in the finite formal certificate.

Research continues on cross-prime unmatched contributions, uniform
moving-parameter bounds, arithmetic coverage, and independent geometric
and incidence constructions. A failed precise certificate only retires
that formulation unless a proof addresses its parent route as well.
